% ==================
% 6) REDEMPTION EXECUTION & FEASIBILITY
% ==================
\section{Redemption execution and feasibility}\label{sec:efficacy}

A position in liquidation exits only through the issuer's own subscription and redemption
process at the published valuation. Any on-chain secondary market is out of scope for this version, and an asset with a bridge or cross-chain dependence sits outside the framework's scope altogether (Section~\ref{sec:overview-channel}). The position is warehoused through the clearing window instead: it is held from
the breach to cash settlement, over the clearing window $\mathrm{LH}^{(a)}$
of Eq.~\eqref{eq:mpor-eff}. The framework reserves
the discount that compensates whoever clears the position, and it does not prescribe the
clearing mechanism. The holder is the deployment or an entity it designates rather than an
outside counterparty standing ready at the moment of need
(Appendix~\ref{sec:free-parameters}).

The discount is bounded by the haircut less its layer total. Liquidation
begins when a position's loan-to-value ratio reaches $\theta^{(a)}$. The
over-collateralisation cushion above the liquidation threshold is
therefore $1-\theta^{(a)} \ge H^{(a)} = \mathrm{IM}^{(a)} +
\delta_{\mathrm{clear}}^{(a)} + \delta_{\mathrm{carry}}^{(a)} + \Lambda^{(a)}$ (equality at the
Prime tier), and whatever discount compensates the clearing party is paid out of it:
\begin{equation}\label{eq:clearing-solvency}
\mathrm{discount} \;\le\; \mathrm{IM}^{(a)} + \delta_{\mathrm{clear}}^{(a)} + \delta_{\mathrm{carry}}^{(a)}
\;=\; H^{(a)} - \Lambda^{(a)} \;\le\; H^{(a)} .
\end{equation}
The bound is non-negative at every admissible
window (Section~\ref{sec:theory:limits}). Whether a party clears at that reserved
level is a stated assumption of the framework (assumption A17 of the assumptions addendum).
The clearing discount is reserved rather than fitted to the valuation history.

\subsection{Redemption capacity}\label{sec:redemption_capacity}

\subsubsection{Issuer capacity and committed settlement}
The redemption capacity of collateral $a$ over one dealing window is
\begin{equation}\label{eq:qcapred}
Q_{\mathrm{cap,red}}^{(a)} \;=\; \min\!\left\{\, Q_{\mathrm{issuer}}^{(a)},\ Q_{\mathrm{queue}}^{(a)} \,\right\},
\qquad
Q_{\mathrm{issuer}}^{(a)} \;=\; R_{\mathrm{cap}}^{(a)},
\end{equation}
where $R_{\mathrm{cap}}^{(a)}$ is the issuer's attested per-window redemption capacity
(USD), obtained from the issuer or transfer agent, used as the one-window
capacity consistently with the caps, the concentration threshold $T$, and the coverage
ratio (Section~\ref{sec:theory:capacity}). The normalisation of attested capacity onto
the dealing window is part of the definition of $R_{\mathrm{cap}}^{(a)}$ there. For an
evergreen fund dealing monthly under a hard quarterly redemption gate the question
inverts structurally. The gate rather than the dealing calendar sets what one window can clear, so the gate window is the governing dealing window, and the window-normalisation rule of Section~\ref{sec:theory:capacity} is met on that basis. $Q_{\mathrm{queue}}^{(a)}$ is specified from the issuer's own documents, the subscription documents together with the offering memorandum's gate and dealing provisions, rather than estimated
from market data.

\subsubsection{The two endpoints of queue capacity}
Feasibility turns on the worst credible recovery in a saturated window, and
that recovery has exactly two endpoints.\footnote{Full derivation. Let $g^{(a)}$ be the
documented gate fraction per dealing window, $A^{(a)}$ the fund AUM, $q_L$ the redemption
request submitted for the position in liquidation, and $q_O^{*}$ the stressed concurrent
third-party demand. The window fills $\min\{q_L+q_O^{*},\,g^{(a)}A^{(a)}\}$, and under
documented allocation the recovery on $q_L$ is
$Q_{\mathrm{queue}}^{(a)} = q_L\cdot\min\{1,\,g^{(a)}A^{(a)}/(q_L+q_O^{*})\}$
(pro-rata) or
$Q_{\mathrm{queue}}^{(a)} = \min\{q_L,\,\max(0,\,g^{(a)}A^{(a)}-q_O^{*})\}$
(first-come-first-served, with $q_L$ at the back of the queue). The stress input
$q_O^{*}$ is prescribed: the largest
attested historical window demand where the transfer agent reports queue depth and
gate utilisation, else the conservative floor
$q_O^{*}=g^{(a)}A^{(a)}$ (a fully subscribed window) where unattested, consistent
with the missing-data rule of Section~\ref{sec:agg-missing}.}
\begin{equation}\label{eq:queue}
Q_{\mathrm{queue}}^{(a)} \;=\;
\begin{cases}
\, 0
  & \parbox[t]{0.5\linewidth}{\raggedright the subscription documents commit to no measurable within-window processing level, so within-window processing is not guaranteed and the conservative saturated-window floor applies,}\\[0.6em]
\, \ge \; q_L^{\mathrm{committed}}
  & \parbox[t]{0.5\linewidth}{\raggedright a documented or attested committed settlement level guaranteeing processing within the window.}
\end{cases}
\end{equation}
The committed-settlement term defaults to the within-window processing level committed in
the issuer's subscription documents. That level is read case by case per listing, since
those terms vary from one issuer to the next. A higher attested committed settlement
level is recognised where one exists.

The borrow cap is zero on the first branch of Eq.~\eqref{eq:queue}, so the
asset can be listed and cannot be borrowed against, and the launch ramp of
Section~\ref{sec:policy-caps} applies only once the committed-settlement term is positive.
On the second branch, $Q_{\mathrm{queue}}^{(a)}$ is bounded
below by the committed settlement level $q_L^{\mathrm{committed}}$, and the check passes
only where that level covers the borrow cap. The same evidence shortens the clearing window
$\mathrm{LH}^{(a)}$ of Eq.~\eqref{eq:mpor-eff}, and neither reduction is a condition of
listing. Documented pro-rata terms do not by themselves establish a guaranteed
within-window clear. They feed the stressed realised $\mathrm{LH}^{(a)}$ and the
clearing-discount reserve (Eq.~\eqref{eq:m-scaling}, Section~\ref{sec:tax-redemption})
rather than a positive committed-settlement term.

\subsubsection{Clearing window per collateral}
The clearing window $\mathrm{LH}^{(a)}$ over which the position is warehoused is the
window of Eq.~\eqref{eq:mpor-eff}, fixed by the three-source hierarchy of
Section~\ref{sec:theory:setting}. It is the same
quantity on which $\delta_{\mathrm{clear}}^{(a)}$ is built (Section~\ref{sec:theory}).

An attested $\mathrm{LH}$, and the $\delta_{\mathrm{clear}}$ and liquidation threshold
derived from it, hold only while the attested terms behind the governing window remain in
force. When they
lapse, the next source in the hierarchy of Eq.~\eqref{eq:mpor-eff} governs and the limits
of Section~\ref{sec:policy} re-derive accordingly.

\subsection{Feasibility condition}\label{sec:feasibility}
The sole capacity check of the framework is that the largest permissible
debt position exits within one dealing window:
\begin{equation}\label{eq:feasibility}
B_{\max}^{(a)} \;\le\; Q_{\mathrm{cap,red}}^{(a)}
\qquad \text{over } \mathrm{LH}^{(a)},
\end{equation}
with $q_L = B_{\max}^{(a)}$ in Eq.~\eqref{eq:queue}. The prescribed stress is the
pool's entire debt at the borrow cap, measured at the point of liquidation, clearing one
window with the whole position redeemed. The interest accruing over that window is reserved
separately by $\delta_{\mathrm{carry}}^{(a)}$. Only the debt amount must clear the window, since the proceeds of the issuer
redemption are applied to the debt first. Any collateral above the repaid debt
queues into the following window at the borrower's risk rather than the
suppliers'. Since Section~\ref{sec:policy} sets
$B_{\max}^{(a)} \le R_{\mathrm{cap}}^{(a)}$ by construction, the issuer capacity term of
Eq.~\eqref{eq:qcapred} holds automatically. The operative content of
Eq.~\eqref{eq:feasibility} is the committed-settlement term, which by Eq.~\eqref{eq:queue} reduces to
the documented committed settlement level.

The check is computed in one pass, in one direction: the caps from Section~\ref{sec:policy} feed
in, Eq.~\eqref{eq:feasibility} either passes or fails, with no iteration
back into the haircut. The check uses the current attested inputs. If the condition fails, limits must tighten:
$B_{\max}^{(a)}$ and $S_{\max}^{(a)}$ are
reduced until Eq.~\eqref{eq:feasibility} passes. Where debt can be attributed to entities, the horizon for the largest single debt is $\lceil B_{(1)}^{(a)}/R_{\mathrm{cap}}^{(a)}\rceil$ windows, and under the supply cap alone a full redemption completes within $M$ windows (Section~\ref{sec:theory:capacity}).

\subsection{Settlement evidence and model consequences}\label{sec:warehouse-ops}
Every parameter in this section is set on the floor state: no measurable committed
settlement level in the subscription documents, no clearing history, and no party standing
ready at the moment of need. Qualifying settlement evidence earns a practitioner-set reduction against that floor (free parameter, Appendix~\ref{sec:free-parameters}). The model recognises three forms of evidence and one live control:
\begin{enumerate}[label=(\roman*),itemsep=1pt]
\item \textbf{The committed settlement level.} The subscription documents of a listed
collateral give the processing level committed within the window.
\item \textbf{Test redemptions.} A completed test corroborates the attested $\mathrm{LH}^{(a)}$. A failed test, or one that settles beyond the attested window, widens it to the documented worst case. The operations companion specifies the test procedure.
\item \textbf{Settlement records.} A realised window longer than the attested one widens $\mathrm{LH}^{(a)}$ and raises the clearing-discount
reserve of Eq.~\eqref{eq:m-scaling}. That record can only tighten the parameters.
\item \textbf{Rate limit on queue saturation.} Queue saturation or issuer processing delays trigger a rate limit on new borrows. The operations companion specifies the release procedure (Section~\ref{sec:operations}).
\end{enumerate}

\subsubsection{Redemption timing evidence}\label{sec:redemption-ops}
Redemption evidence compares the realised time from notice to settlement with the attested window. An overrun widens $\mathrm{LH}^{(a)}$ under item~(iii) above. The operations companion specifies the notice, reconciliation, and recordkeeping procedures (Section~\ref{sec:operations}).
